\section{Static arbitrage in total variance}
\label{app:arbitrage}

Fix a maturity $T$ with forward $F$ and discount factor $D$, and write the call price as
$C(K)=DF\,c(k)$ with $k=\log(K/F)$ and $c(k)=c_{\text{B}}\bigl(k,w(k)\bigr)$, where
\[
  c_{\text{B}}(k,w)=N(d_1)-e^{k}N(d_2),\qquad d_{1,2}=-\frac{k}{\sqrt w}\pm\frac{\sqrt w}{2}.
\]
Black's partial derivatives are $\partial_k c_{\text{B}}=-e^kN(d_2)$,
$\partial_w c_{\text{B}}=e^k\phi(d_2)/(2\sqrt w)>0$ and
$\partial_{kk}c_{\text{B}}-\partial_k c_{\text{B}}=e^k\phi(d_2)/\sqrt w=2\,\partial_w c_{\text{B}}$, using
$\phi(d_1)=e^k\phi(d_2)$.

\paragraph{Calendar.} With deterministic rates and proportional dividends, a calendar spread at a
fixed forward-moneyness is free of arbitrage if and only if $c(k,T)$ is nondecreasing in $T$ at fixed
$k$~\citep{carr2005,gatheral2014}. Since $\partial_w c_{\text{B}}>0$, this holds if and only if
$\partial_T w(k,T)\ge 0$.

\paragraph{Butterfly.} The risk-neutral density of the underlying at $T$ is
$q(K)=D^{-1}\partial^2_K C$~\citep{breeden1978}. Since $\partial_K=K^{-1}\partial_k$,
$\partial^2_K C=DF K^{-2}(c''-c')$, with primes the total derivatives in $k$. Expanding $c'$ and $c''$
through $w(k)$,
\[
  c''-c'=(\partial_{kk}-\partial_k)c_{\text{B}}+(2\partial_{kw}-\partial_w)c_{\text{B}}\,w'
  +\partial_{ww}c_{\text{B}}\,w'^2+\partial_w c_{\text{B}}\,w''.
\]
Differentiating $\partial_w c_{\text{B}}$ gives $\partial_{kw}c_{\text{B}}=\partial_w c_{\text{B}}\,(1+d_2/\sqrt w)$ and
$\partial_{ww}c_{\text{B}}=\partial_w c_{\text{B}}\,\bigl(k^2/(2w^2)-1/(2w)-1/8\bigr)$. Substituting
$d_2/\sqrt w=-k/w-1/2$ and collecting terms,
\[
  c''-c'=2\,\partial_w c_{\text{B}}\Bigl[1-\frac{kw'}{w}+\Bigl(\frac{k^2}{4w^2}-\frac{1}{4w}-\frac1{16}\Bigr)w'^2+\frac{w''}{2}\Bigr]
  =2\,\partial_w c_{\text{B}}\;g(k),
\]
with $g$ as in \cref{eq:conditions}. Hence $q(K)=\phi\bigl(d_2(k)\bigr)\,g(k)/\bigl(K\sqrt{w(k)}\bigr)$,
which is nonnegative exactly where $g$ is~\citep{gatheral2014}.

\paragraph{Wings.} For any arbitrage-free smile, $\limsup_{|k|\to\infty}w(k)/|k|\le2$~\citep{lee2004}.
The penalized network enforces it as a bound on the slope of $w$ at the outermost collocation
strikes; eSSVI satisfies it through $\max(u_+,u_-)<4$.

\paragraph{Grids.} Violation shares are evaluated on 40 maturities and 101 strikes per maturity (801
for the calendar check) spanning the quoted domain, and on the same grid widened by 50\% of the
quoted width on each side. A calendar violation counts when $w$ falls between consecutive grid
maturities by more than a numerical tolerance; a butterfly violation when $g<-10^{-5}$.

\section{eSSVI parametrization and price mixtures}
\label{app:essvi}

\paragraph{Feasible knots.} At knot maturities $T_1<\dots<T_n$, unconstrained rates
$(a_i,b_i^+,b_i^-)$ map to knots as follows, with $\operatorname{sp}$ the softplus,
$\Delta_i=T_i-T_{i-1}$ and $\delta_i=\log T_i-\log T_{i-1}$:
\[
  u^\pm_i=4\Bigl(1-\exp\Bigl(-\sum_{l\le i}\operatorname{sp}(b^\pm_l)\,\delta_l\Bigr)\Bigr),\qquad
  \theta_i=\max\Bigl(\theta_{i-1}\frac{\psi_i}{\psi_{i-1}},\;\frac{\psi_i\max(u^+_i,u^-_i)}{4}\Bigr)+\operatorname{sp}(a_i)\,\Delta_i,
\]
with $\psi_i=(u^+_i+u^-_i)/2$. The wings are nondecreasing and below 4. Because $\psi$ is
nondecreasing, $\theta_i>\theta_{i-1}$; the first term of the maximum keeps $\psi/\theta$
nonincreasing and the second is the butterfly bound $\psi\max(u_+,u_-)\le4\theta$. Every condition
of \cref{sec:models} therefore holds for every rate vector, and every feasible knot sequence with
strict increases has a preimage.

\begin{proposition}
\label{prop:mixture}
Let $c_1,\dots,c_M$ be normalized call price surfaces on a common set of expiries, each free of static
arbitrage, and let $\pi_j\ge0$ with $\sum_j\pi_j=1$ be constant in $k$ and $T$. Then
$c=\sum_j\pi_jc_j$ is free of static arbitrage, and its Black implied total variance $w(k,T)$ exists,
is unique, lies between $\min_j w_j$ and $\max_j w_j$, and satisfies \cref{eq:conditions}.
\end{proposition}

\begin{proof}
For each $T$, every $c_j$ is convex and nonincreasing in the strike and lies between
$(1-e^k)^+$ and 1; for each $k$, it is nondecreasing in $T$. Each of these properties is preserved
by a convex combination with weights that do not depend on $k$ or $T$, so $c$ has them too. Since
$w\mapsto c_{\text{B}}(k,w)$ is a strictly increasing bijection from $(0,\infty)$ onto
$\bigl((1-e^k)^+,1\bigr)$, the implied variance exists and is unique, and because $c$ lies between
$\min_j c_j$ and $\max_j c_j$ it lies between the components' variances. Monotonicity of $c$ in $T$
gives $\partial_Tw\ge0$ (\cref{app:arbitrage}). The density of $c$ is $\sum_j\pi_jq_j\ge0$, so
$g\ge0$.
\end{proof}

If the weights varied with $T$, a component with low variance could gain weight at a longer maturity
and the mixture price could fall, which is why \mixfour{} uses one set of weights per day.

\section{Forwards and discount factors from put--call parity}
\label{app:forwards}

For each expiry, take the 20 strikes nearest the forward (smallest $|C-P|$) at which both the call
and the put are quoted two-sided. Put--call parity, $C-P=DF-DK$, makes the mid price difference
linear in $K$ with slope $-D$ and intercept $DF$. The slope is weakly identified at short maturities,
where $D\approx1-rT$, so estimation runs in three stages per date: (i) a weighted regression per
expiry with $T\ge30$ days gives $r=-\log D/T$; (ii) a robust kernel smoother in $\log T$ turns these
into a rate curve, flat below 30 days; (iii) with $D=e^{-r(T)T}$ for every expiry, $F$ is the
weighted median of $K+(C-P)/D$ over the pairs. An expiry is valid when the relative median absolute
deviation of those per-strike forwards is below 0.002.

\section{Robust refits}
\label{app:robust}

The least-squares objective of a surface is $\frac1n\sum_i u_i^2$ with
$u_i=(\sigma^{\text{model}}_i-\sigma^{\text{mkt}}_i)/s_i$, where $s_i$ is the IV half-spread with a floor
of 0.1\vp. One step of iteratively reweighted least squares for the Cauchy loss
$\rho(u)=\tfrac{c^2}{2}\log\bigl(1+u^2/c^2\bigr)$ weights each quote by
$\rho'(u_i)/u_i=1/\bigl(1+u_i^2/c^2\bigr)$, the residual from the previous pass. This is the squared
loss with each half-spread replaced by $s_i\sqrt{1+u_i^2/c^2}$, so every refit reuses the
least-squares solver. The pooled variant replaces $|u_i|$ in the weight by
$\max\bigl(|u_i|,|\bar u_i|\bigr)$, where $\bar u_i$ is the mean residual of quote $i$ and up to two
strikes on either side in its expiry. Quotes that are legs of model-free arbitrage are scored but
excluded from every fit.

\section{Reversion tests}
\label{app:reversion}

For origin day $t$ and horizon $h$ (in screened trading days), let $r_{i,t}$ be the residual of
contract $i$ in price half-spreads and keep the contracts quoted on both $t$ and $t+h$. For flagged and
unflagged contracts separately, the slope $b_t$ of the cross-sectional regression
$r_{i,t+h}-r_{i,t}=a_t+b_t r_{i,t}+e_{i,t}$ is estimated on each day with at least five flags. The
reported slopes are the means of $b_t$, and the test of a difference uses the daily series
$b^{\text{flag}}_t-b^{\text{unflag}}_t$ with Newey--West standard errors~\citep{newey1987}, in the
manner of \citet{fama1973}; pooled regressions over thousands of quotes a day would overstate the
significance.

\paragraph{Decomposition.} With $x_{i,t}=(\sigma^{\text{mkt}}_{i,t}-\sigma^{\text{model}}_{i,t})/s_{i,t}$,
the change $x_{i,t+h}-x_{i,t}$ is not exactly the change of the residual, because $s$ changes between
days. Measured in the origin day's half-spread,
\[
  \frac{(\sigma^{\text{mkt}}_{i,t+h}-\sigma^{\text{model}}_{i,t+h})-(\sigma^{\text{mkt}}_{i,t}-\sigma^{\text{model}}_{i,t})}{s_{i,t}}
  =\underbrace{\frac{\sigma^{\text{mkt}}_{i,t+h}-\sigma^{\text{mkt}}_{i,t}}{s_{i,t}}}_{\text{market}}
  -\underbrace{\frac{\sigma^{\text{model}}_{i,t+h}-\sigma^{\text{model}}_{i,t}}{s_{i,t}}}_{\text{model}},
\]
and since an OLS slope is linear in the dependent variable, the slope of the left side on $x_{i,t}$ is
the market slope minus the model slope, day by day and on average. A negative market slope means
prices move towards the surface. The total in this basis differs from the price-residual slope; on
SPX 2022 the robust detector's one-day flagged slope is $-0.46$ in prices, $-0.33$ in IV and $-0.21$
in IV with the origin day's spread, so each part is compared with its own total.

\paragraph{Forecast tests.} The Diebold--Mariano statistic~\citep{diebold1995} is the mean of the
daily loss differential (baseline minus model) over its Newey--West standard error, with Bartlett
bandwidth $\max\bigl(h,\,4(n/100)^{2/9}\bigr)$ because overlapping $h$-day errors are autocorrelated
up to lag $h-1$, and a two-sided $p$-value from $t_{n-1}$.

\section{Backtest accounting}
\label{app:backtest}

A signal on day $s$ (sell if rich, buy if cheap) is executed at the close of day $s+\ell$, with
$\ell=1$ (next day) or $\ell=0$ (same day), in $n=V/(100\cdot0.01\,\nu)$ contracts, where $V=\$1{,}000$ is the
target dollar vega per vol point, 100 the contract multiplier and $\nu$ the Black vega per unit of
volatility at the signal. The
entry and the exit each cost $100\,n$ half-spreads plus \$1 per contract. Between closes $t$ and
$t+1$ a position of signed size $q=\pm100n$ earns
\[
  q\bigl[(m_{t+1}-m_t)-\Delta_t\,(F_{t+1}-F_t)\bigr],
\]
with $m$ the mid price, $\Delta_t$ the Black forward delta at the day's market IV and $F$ the expiry's
parity forward; the hedge is held in the forward. The position is closed $\ell$ days after the first
close at which its residual is below half a spread, or after five days. The panel holds
out-of-the-money options only, so an option that moves into the money is marked through parity from
the out-of-the-money option at its strike. A contract that drops out of the cleaned panel keeps its
hedge for up to three days and is marked when it is quoted again; beyond that it is closed at its
last quote. Entries are restricted to contracts that stay inside the panel until their latest
possible exit, and a contract is traded at most once at a time.

Each trade has 20 placebos, each a random unflagged and tradable contract from the signal day with
the same type, direction and maturity--moneyness bucket, entered and exited on the same days. The
strategy is compared with the mean placebo by a Newey--West test on the daily difference of net PnL.
Sharpe ratios are the mean of daily net PnL over its standard deviation, times $\sqrt{252}$.

\section{Forecasts}
\label{app:forecast}

\begin{table}[h]
  \centering\small
  \caption{Forecast RMSE (vp) on every quote of the target day, $h$ trading days ahead, on the
  validation years. RW is a random walk on the latent or on the origin day's \mixfour{} fit. The $p$-value is
  a Diebold--Mariano test~\citep{diebold1995}, with Newey--West variance~\citep{newey1987}, of the
  VAR against the \mixfour{} random walk; the last column scores the target day's own decoded surface, a floor
  for any forecast through the decoder.}
  \label{tab:forecast}
  \begin{tabular}{llrrrrr}
\toprule
panel & $h$ (days) & VAR(1) & latent RW & \texttt{mix4} RW & $p$ & decoded target \\
\midrule
SPX 2022 & $1$ & $1.73$ & $1.73$ & $\mathbf{1.35}$ & $<0.001$ & $1.19$ \\
 & $5$ & $2.70$ & $2.70$ & $\mathbf{2.53}$ & $0.15$ & $1.20$ \\
 & $20$ & $\mathbf{3.90}$ & $4.25$ & $4.27$ & $0.27$ & $1.22$ \\
Heston 2021 & $1$ & $0.64$ & $\mathbf{0.60}$ & $0.60$ & $<0.001$ & $0.20$ \\
 & $5$ & $1.30$ & $\mathbf{1.13}$ & $1.15$ & $0.03$ & $0.20$ \\
 & $20$ & $2.82$ & $\mathbf{1.93}$ & $1.96$ & $0.01$ & $0.20$ \\
\bottomrule
\end{tabular}

\end{table}

\section{The prediction written before the last test run}
\label{app:prediction}

The rule of \cref{sec:refits} was frozen on the validation years and its configurations committed
(\texttt{30c766c}) on November 2, 2026, before any test panel was refitted. The validation notes
written then, which are kept with the results rather than in the repository, ended with this
prediction for the test run, quoted verbatim:

\begin{quote}
\small\ttfamily\raggedright
PREDICTION for test: next-day residual "reversion" of flags is mostly not the market moving;
flagged market-side move not distinguishable from unflagged.
\end{quote}

\noindent The same notes recorded the validation results that motivated it: on the SPX 2022 refits
the robust detector's flags still reverted faster than unflagged ones at one day in price residuals
(difference $-0.27$), but their market part was small and not significant ($-0.10$, $p=0.62$), and
the flagged market move was not distinguishable from the unflagged one ($p=0.69$).
